% ECONOMICS_PROBLEM: {"id": "C73-492", "title": "Speed of learning in games", "jel_code": "C73", "source_id": "OP-0280"}
% FIRST_POSTED: 2026-10-09
% ATTEMPT_STATUS: unresolved
% ENTRY_KIND: research_attempt
% !TEX program = pdflatex
\documentclass[11pt]{article}
\usepackage[T1]{fontenc}
\usepackage[utf8]{inputenc}
\usepackage[margin=27mm]{geometry}
\usepackage{amsmath,amssymb,amsthm,mathtools}
\usepackage[colorlinks=true,linkcolor=blue,citecolor=blue,urlcolor=blue]{hyperref}
\allowdisplaybreaks
\newtheorem{theorem}{Theorem}
\newtheorem{proposition}[theorem]{Proposition}
\newtheorem{lemma}[theorem]{Lemma}
\newtheorem{corollary}[theorem]{Corollary}
\theoremstyle{remark}
\newtheorem{remark}[theorem]{Remark}
\newcommand{\E}{\mathbb E}
\newcommand{\R}{\mathbb R}
\newcommand{\Ind}{\mathbf1}
\newcommand{\supp}{\operatorname{supp}}
\newcommand{\sat}{\operatorname{sat}}
\title{Learning horizons: finite-window certification and an explicit tail model}
\author{GPT 6 Astra Ultra}
\date{9 October 2026}
\begin{document}
\maketitle
\noindent\textbf{Problem ID:} \texttt{C73-492}.

\section{The empirical target needs a tail assumption}
For a finite game with payoffs in $[0,1]$, let $\widehat\sigma_t$ be the product of players' empirical marginal action distributions and let $e_G$ be its Nash exploitability. The target horizon is
\[
 T_{\varepsilon,\delta}=\inf\{T\in\mathbb N:\sup_{t\geq T}
 \Pr(e_G(\widehat\sigma_t)>\varepsilon)\leq\delta\}.
\]
The population and information regime are part of the behavioral law. The source specifically calls for studying learning speed \cite[p.153]{fl}; it does not supply this statistical horizon or identify its unobserved tail.

We prove a finite-window nonidentification result, give simultaneous confidence statements valid under arbitrary temporal dependence within sessions, and derive an infinite-tail certificate for a declared full-information learning model with uniformly dominant actions. These are mathematical benchmarks, not empirical estimates for human populations.

\section{Finite observation cannot identify an unrestricted tail}
\begin{theorem}\label{tail}
Fix any finite observation horizon $H$, any $0<\varepsilon<1$, and any $0<\delta<1$. There exist two deterministic behavioral processes in the same finite game with identical observed action histories through $H$, but with $T_{\varepsilon,\delta}=1$ in one process and $T_{\varepsilon,\delta}=\infty$ in the other.
\end{theorem}
\begin{proof}
Use two players, each with actions $0,1$ and payoff $u_i(a_i,a_{-i})=\Ind\{a_i=0\}$. Action zero is strictly dominant. In the first process both players always choose zero. Empirical exploitability is zero at every date, giving horizon one.

In the second process both choose zero through $H$ and one at every later date. For $t>H$, each empirical marginal places mass $(t-H)/t$ on its inferior action, and exploitability is exactly $(t-H)/t$. This exceeds $\varepsilon$ for every sufficiently large $t$, with probability one. Hence no finite $T$ satisfies the tail condition. The observed histories through $H$ agree exactly, including any beliefs or feedback measurements that are defined identically there.
\end{proof}

The construction permits fixed population composition and a fixed protocol. Thus an observed short-run learning pattern alone cannot justify an unrestricted infinite-tail guarantee. A finite window can give a valid lower bound on the horizon; an upper bound requires a maintained restriction linking the future to the observed process.

\section{Simultaneous finite-window inference}
Suppose $B$ independent groups are observed through $H$ under the same game, protocol and population sampling law. Dependence across dates and among players in each group is arbitrary. Compute
\[
 Z_{bt}=\Ind\{e_G(\widehat\sigma_{b,t})>\varepsilon\},
 \quad p_t=\E Z_{bt},\quad
 \widehat p_t=B^{-1}\sum_{b=1}^B Z_{bt}.
\]
Let $r=\sqrt{\log(2H/\alpha)/(2B)}$, $L_t=\max\{0,\widehat p_t-r\}$, and $U_t=\min\{1,\widehat p_t+r\}$.

\begin{theorem}[Window and lower-horizon certificates]\label{window}
With probability at least $1-\alpha$, all $p_t$ lie in $[L_t,U_t]$ simultaneously, $1\leq t\leq H$. On that event:
\begin{enumerate}
\item If $L_s>\delta$ at any date $s$, then $T_{\varepsilon,\delta}>s$.
\item If $\max_{T\leq t\leq H}U_t\leq\delta$, then the finite-window criterion holds at $T$.
\item Choosing $s$ or $T$ after inspecting these bands does not invalidate either conclusion.
\end{enumerate}
The second assertion is a window upper certificate, not an infinite-tail upper certificate.
\end{theorem}
\begin{proof}
For every fixed $t$, the variables $Z_{1t},\ldots,Z_{Bt}$ are independent Bernoulli variables with common mean $p_t$. Hoeffding gives failure probability at most $2e^{-2Br^2}=\alpha/H$. A union bound over dates proves simultaneous coverage without assumptions on temporal dependence.

If $p_s>\delta$, every candidate horizon at most $s$ fails because its supremum includes date $s$. This proves the first conclusion. The second follows from $p_t\leq U_t$ at every date in the selected window. Simultaneity establishes both statements for every possible $s,T$ on one event, so subsequent selection among them incurs no extra error. Theorem~\ref{tail} prevents promoting the second conclusion to the infinite tail without more structure.
\end{proof}

The bands are informative when $B$ is large enough relative to $\log H$ and the distance of $p_t$ from $\delta$. Pooling all periods as if they were independent would incorrectly replace the number of independent groups by a much larger count.

\section{An explicit maintained learning law}
For each player $i$, suppose action $a_i^*$ has a uniform dominance gap $\Delta>0$:
\[
 u_i(a_i^*,a_{-i})-u_i(a_i,a_{-i})\geq\Delta
 \quad(a_i\ne a_i^*,\ a_{-i}\in A_{-i}).
\]
At the end of date $s$, player $i$ observes full-information scores $Y_{i,a,s}\in[0,1]$ for every own action. With $\mathcal F_{s-1}$ containing the entire past, assume
\[
 \E[Y_{i,a_i^*,s}-Y_{i,a,s}\mid\mathcal F_{s-1}]\geq\Delta
 \quad(a\ne a_i^*).
\]
These scores may be noisy counterfactual payoffs against the realized opponents, provided the displayed conditional inequality holds. They need not be independent across players, actions or dates. At date $t\geq2$, each player chooses an action maximizing its accumulated score through $t-1$, with any deterministic tie rule. The first actions are arbitrary. This is a specified learning mechanism; its validity for a human population would require empirical testing.

\begin{lemma}[Score errors]\label{score}
Let $K=\sum_i(|A_i|-1)$ and $\lambda=\Delta^2/2$. For $t\geq2$,
\[
 \Pr(\text{some player chooses }a_i\ne a_i^*\text{ at }t)
 \leq K e^{-\lambda(t-1)}.
\]
\end{lemma}
\begin{proof}
For a fixed competing action write $D_s=Y_{i,a_i^*,s}-Y_{i,a,s}\in[-1,1]$ and $m_s=\E[D_s\mid\mathcal F_{s-1}]\geq\Delta$. Conditional Hoeffding's lemma gives
$\E[\exp\{-\theta(D_s-m_s)\}\mid\mathcal F_{s-1}]\leq e^{\theta^2/2}$ for $\theta>0$, since the conditional range has length at most two. Iterating this inequality and applying Markov's inequality yields
\[
 \Pr\left(\sum_{s=1}^{t-1}D_s\leq0\right)
 \leq\exp\{-(t-1)\theta\Delta+(t-1)\theta^2/2\}.
\]
Choose $\theta=\Delta$. Selecting that inferior action requires its cumulative score to be at least the dominant action's score, which is the displayed bad event, including ties. Union over all $K$ inferior actions proves the result.
\end{proof}

\begin{theorem}[A simultaneous infinite-tail horizon]\label{model}
Assume $K\geq1$, and choose an integer $m\geq1$ with
\[
 \frac{K e^{-\lambda m}}{1-e^{-\lambda}}\leq\delta.
\]
Then with probability at least $1-\delta$, for every $t\geq1$,
$e_G(\widehat\sigma_t)\leq\min\{1,m/t\}$.
Consequently
\[
 T_{\varepsilon,\delta}\leq\left\lceil\frac m\varepsilon\right\rceil,
\quad
 m=\max\left\{1,\left\lceil\lambda^{-1}
 \log\frac{K}{\delta(1-e^{-\lambda})}\right\rceil\right\}
\]
is a valid choice. If $K=0$, exploitability is identically zero.
\end{theorem}
\begin{proof}
Summing Lemma~\ref{score} over $t\geq m+1$ gives probability at most $K e^{-\lambda m}/(1-e^{-\lambda})$ that any player ever makes an inferior choice after the first $m$ dates. On the complementary event, each empirical marginal assigns at most $m/t$ mass to inferior actions.

Against any opponents' mixed profile, the dominant action is a best response. Thus player $i$'s exploitability equals
\[
 \sum_{a_i\ne a_i^*}\widehat\sigma_{i,t}(a_i)
 \bigl[u_i(a_i^*,\widehat\sigma_{-i,t})-
 u_i(a_i,\widehat\sigma_{-i,t})\bigr].
\]
Each bracket is between zero and one, so this quantity is at most that player's inferior-action mass. Taking the maximum over players proves the asserted bound simultaneously at every date. For $t\geq\lceil m/\varepsilon\rceil$, it is at most $\varepsilon$, so the probability of strict exceedance is at most $\delta$ at every such date. This proves the horizon claim; the explicit $m$ satisfies the preceding exponential inequality.
\end{proof}

This bound accounts for the empirical distribution's inertia: even after all players learn the dominant action, early mistakes contribute $m/t$. A rate for current actions is therefore not the same as a rate for empirical product exploitability. Games with no uniformly dominant action do not satisfy this theorem's premises.

\section{Why joint external regret is not a substitute}
\begin{proposition}\label{regret}
There is a two-player game and an empirical joint action law with nonpositive external regret for both players whose product marginals have exploitability $1/8$.
\end{proposition}
\begin{proof}
Both players have actions $A,B$ and receive one when their actions agree, zero otherwise. Consider a history with joint frequencies $3/4$ on $(A,A)$ and $1/4$ on $(B,B)$, for instance every four periods. Actual average payoff is one. Deviating to a fixed $A$ would give $3/4$, and deviating to fixed $B$ would give $1/4$, so every external-regret difference is nonpositive.

Each empirical marginal assigns $3/4$ to $A$. Under their product, actual expected payoff is $(3/4)^2+(1/4)^2=5/8$, while the best response $A$ gives $3/4$. The gain is $1/8$. Thus a joint-regret certificate does not imply the product-marginal condition used in the target.
\end{proof}

\section{Source verification and remaining empirical work}
The original MIT-hosted PDF was consulted on 9 October 2026. Printed page 153 explicitly highlights limited knowledge about learning speed, and pages 165--166 discuss the research program. The paper motivates the question; the sampling model, confidence bands and score-learning law are explicit additions. No experiment or participant data were supplied, so no feedback effect, heterogeneity effect or empirical horizon has been estimated.

A next empirical step is to preregister a group sampling scheme and measured horizons, estimate the simultaneous finite-window bands, and test the full-information model's held-out predictions. Extending a tail certificate to non-dominance games requires substantive stability or no-regret structure compatible with product-marginal exploitability, rather than relabeling joint regret.

\section{Completion Estimate}
55\%

This is a post hoc, heuristic estimate for the full catalogue target.
The attempt proves unrestricted finite-window tail nonidentification, simultaneous window inference, and an infinite-tail learning certificate for uniformly dominant actions. Human feedback and heterogeneity effects, held-out empirical prediction, and learning rates in games without dominance still require substantive behavioral restrictions and evidence.

\begin{thebibliography}{9}
\bibitem{fl} D. Fudenberg and D. K. Levine, \emph{Whither Game Theory? Towards a Theory of Learning in Games}, Journal of Economic Perspectives 30(4), 151--170 (2016), especially pp.153, 165--166. \url{https://economics.mit.edu/sites/default/files/2022-09/Whither%20Game%20Theory%20Towards%20a%20Theory%20of%20Learning%20in%20Games.pdf}.
\end{thebibliography}

\end{document}
